\begin{proposition}[The dilatation and the translations on the families from the defect cycles]
\label{prop:c11-masterfam}
With the master formula \eqref{eq:c11-decompx} along $\Delta$ and $P_n$, the sector coefficients $-\tfrac13$ of Theorem \ref{thm:c9-theta2} and the families and kinematics of Definition \ref{def:c11-families}, on the parts with the fewest fields:
\begin{itemize}
\item (i) on the $\mathfrak{su}(2)$ sector, $L\le8$, $D_1=-8\,\text{units}\cdot\sum_l(1-P_{l,l+1})$, $P_{l,l+1}$ the exchange of the letters at the sites $l$, $l+1$, which is $\frac{g^2N}{4\pi^2}$ times the Heisenberg Hamiltonian; on the $\mathfrak{su}(1|1)$ sector, $L\le6$, on all twenty-two graded cyclic sectors, $D_1=-8\,\text{units}\cdot\sum_l(1-\Pi_{l,l+1})$ with the graded exchange $\Pi_{l,l+1}$;
\item (ii) on twist two, $S'\le6$, the spectrum of $D_1$ is $\{-32\,h(j)\ \text{units}: j\le S'\ \text{even}\}$, and on the fermion $\mathfrak{sl}(2)$ sector, $S'\le5$, it is $\{-32\,h(j+1)\ \text{units}: j\le S'\ \text{odd}\}$;
\item (iii) on the family of two gaugino components with $Z'F_{11}$, $S'\le3$, the eigenvalues are $0$, $-48$, $-\tfrac{200}3$ units, on the gluon $\mathfrak{sl}(2)$ sector, $S'\le2$, $-48$, $-\tfrac{200}3$ units, and on $(F_{11},\psi)$, $S'\le4$, $-48$, $-\tfrac{200}3$, $-\tfrac{392}5$ units: $-8$ units times the eigenvalues of Proposition \ref{prop:c11-targets};
\item (iv) the sum over the pairs of $\Upsilon_{P_n}$ is $-\tfrac13$ units times $\partial_n$ on twist two and on the fermion $\mathfrak{sl}(2)$ sector, and on the families of (iii), and $D_1(\partial_nO)=\partial_nD_1(O)$ on twist two and the families of (ii), (iii), as an identity of matrices from $S'$ to $S'+1$;
\item (v) pair by pair, $\Upsilon_\Delta(X_a,X_b)=4\,\mathrm{lc}(X_aX_b)+\big(\text{composite pair kernel along }\Delta\big)$, $\Upsilon_{P_b}$ and $\Upsilon_{M_\omega}$ equal to the composite pair kernels along $P_b$ and $M_\omega$, and $\Upsilon_R=0$, on every pair of slots with defect cycles and every leg set with at most one letter derivative on which both evaluations are complete, and with at most two letter derivatives in the simplex evaluation.
\end{itemize}
The spectra of (i)--(iii) are those of Table \ref{tab:c11-results}, obtained here without the composite term.
\end{proposition}

\begin{proof}
The two-body kernels are the coefficients of $\varepsilon^0$ of the defect cycles of the pairs of letters with $S^{(1)}(\Delta)$ or $S^{(1)}(P_n)$ at the special edge or on a line, evaluated exactly by Theorem \ref{thm:c11-oddsimplex} with the colour words reduced to the planar adjacent form of Lemma \ref{lem:c10-colour} (iv), and the one-body terms are the brackets with \eqref{eq:c9-F1Delta} and with $\hbar F_1^{(1)}(P_n)$; the operators are diagonalised exactly on bases of the families closed under them, with the projection to cohomology of Remark \ref{rem:c11-method}. Items (i), (ii), (iv) on twist two and the fermion sector and (v) at one letter derivative are obtained in both evaluations of Remark \ref{rem:c11-oddmethod}; the families of (iii), the matrix identity of (iv) on them and (v) at two letter derivatives in the simplex evaluation. Item (v) summed over the pairs of a cocycle is Corollary \ref{cor:c11-DPdefect}.
\end{proof}

\begin{corollary}[The components on the scalar bilinears: representatives and invariant content]
\label{cor:c11-F0}
Let $O=\int\omega\,(\iota_{w_1}A,\iota_{w_2}A)$, $\psi=0$, $L=\log\tfrac32$, $\mathcal B_1,\dots,\mathcal B_4$ as in Lemma \ref{lem:c11-cocycle}, and let the action-side solutions be those of Theorem \ref{thm:c9-quadratic} (iii), Theorem \ref{thm:c9-P1} (i) and Proposition \ref{prop:c11-oddsol}, with the uniform coefficients $t_\Delta$, $t_P$, $t_M$, $t_R$, $t_{\xi_\pm}$, $t_{\eta_\pm}$ of Theorem \ref{thm:c11-oddar2}. The representatives of the components given by the master formula on the part with two legs $A|_{W_{\mathbb R}^\perp}$ are, in units,
\[
\begin{aligned}
&Q_1(O)=\big(\tfrac{79}{12}-3L-t_\xi\big)\mathcal B_1+4L\,\mathcal B_2+4L\,\mathcal B_3+0\cdot\mathcal B_4,\qquad \hbar L_1(O)^{(0)}=-(2+4L)\,(w_1,w_2)\int\omega\,(dc,A|_V),\\
&S_1(O)=\big(\tfrac{79}{12}-3L-t_\eta\big)\{S^{(1)}(\eta),O\}+x_O^\mu\big(4L\,\mathcal B_2+4L\,\mathcal B_3\big)\big(\rho(e_\mu)\eta\big),\\
&D_1(O)=-\big(4t_\Delta+\tfrac43\big)\,O\ \text{(traceless)},\qquad D_1(O)=-\big(4t_\Delta+\tfrac{148}3\big)\,O\ \text{(trace)},\quad\text{at the point }x_O\text{ of the density},\\
&P_1(O)=-\big(\tfrac13+t_P\big)\,\partial_bO,\qquad M_1(O)=-\big(\tfrac13+t_M\big)\,\{S^{(1)}(M_\omega),O\},\qquad R_1(O)=0,\\
&K_1(O)=-\big(\tfrac13+\theta^K_\perp\big)\{S^{(1)}(K_b),O\}+\begin{cases}0&\text{(traceless)},\\ 48\,(b,x)\,O&\text{(trace)},\end{cases}\ -\ \sum_{j=1}^4\kappa^K_j\big\{G^K_j,O\big\},
\end{aligned}
\]
with $t_\xi$, $t_\eta$ the coefficients of the chirality of $\xi$, $\eta$, the spinor $\rho(e_\mu)\eta$ in $\mathcal B_2$, $\mathcal B_3$ in place of $\xi$, and the intrinsic parts of $S_1$ and $K_1$ equal to $0$. Hence:
\begin{itemize}
\item (i) the scheme-independent content is $D_1=0$ on the twenty traceless bilinears and $D_1=-48$ units times $O$ on the trace, $\frac{3g^2N}{2\pi^2}O$ for $\mathfrak g=\mathfrak{su}(N)$, twice the Konishi anomalous dimension: these are the eigenvalues of the induced operator on the cohomology of $\Lambda^{(0)}$, at the value $t_\Delta=-\tfrac13$ fixed by the arity-two equations, invariant under the inner and outer directions and under the finite parts $\psi'$ (Theorem \ref{thm:c8-invariants} (iii));
\item (ii) as functions of the action-side free directions at $\psi=0$, with $c_Q$, $c_S$ the coefficients of $\mathcal B_1$ and of $\{S^{(1)}(\eta),O\}$ and $p_1=-(\tfrac13+t_P)$, on the traceless bilinears
\[
c_Q(\xi_+)+c_Q(\xi_-)-p_1=c_Q(\xi_+)+c_S(\eta_+)-\tfrac14D_1=\tfrac{27}4-8\log\tfrac32,
\]
independent of $t_\Delta$, $t_P$, $t_\xi$, $t_\eta$ and of the inner and outer directions by \eqref{eq:c11-trel}; a finite part $\psi'$ of the operators shifts $c_Q$, $c_S$ and $p_1$ by the coboundaries of Proposition \ref{prop:c11-freedom} (iv), with weight-dependent coefficients on the levels of Weyl weight $2$, $\tfrac52$ and $3$, so these combinations become scheme-independent together with the one-loop components on the weight-$\tfrac52$ level that enter (R2) and (R3) of Proposition \ref{prop:c11-relations};
\item (iii) $P_1$, $M_1$ and $R_1$ vanish at the values $t_P=t_M=t_\Delta=-\tfrac13$ of the translation-covariant scheme, and $L_1^{(0)}(O)$ is removed by the finite part $\psi(O)=-(1+2L)$ units times $(w_1,w_2)\int\omega\,(A|_V,A|_V)$ (Theorem \ref{thm:c11-bilinears} (iii));
\item (iv) the scheme artefacts are $c_Q$ and $c_S$ separately, every coefficient with $\log\tfrac32$, $P_1$, $L_1^{(0)}(O)$, $M_1$, $R_1$, the one-body part of $K_1$, which depends on $\theta^K_\perp$ and the four directions $\kappa^K_j$ left free on two legs, and the individual two-body kernels.
\end{itemize}
\end{corollary}

\begin{proof}
The two-body parts are Theorem \ref{thm:c11-bilinears} (i), Proposition \ref{prop:c11-SKbil} and Corollary \ref{cor:c11-DPdefect}: $\Upsilon^0_\Delta=\tfrac{40}3O$, $-\tfrac{104}3O$, $\Upsilon_{P_b}=-\tfrac13\partial_bO$, $\Upsilon^0_{M_\omega}=0$, $\Upsilon_R=0$, the intrinsic parts along $\eta$ and $K_b$ vanishing on the colour singlet. The one-body parts are the brackets of $O$ with the solutions; on $O$ only the terms with the leg $A^*|_{W_{\mathbb R}^\perp}$ act. Along $\xi$, $T_2$ is the only such term of $\hbar F_1^{(1)}(\xi)$ and $\{T_2,O\}=\mathcal B_1$, which gives $-(L-\tfrac{79}{12}+t_\xi)\mathcal B_1$; along $\eta$ the position part of the solution is $x^\mu\hbar F_1^{(1)}(\rho(e_\mu)\eta)$ with the uniform $t_\eta$, and $s_0$, $u_0$ have no leg $A^*$; along $\Delta$, $-\{\tfrac{22}3\int(A,A^*),O\}=-\tfrac{44}3O$ and $-\{t_\Delta S^{(1)}(\Delta)_\perp,O\}=-4t_\Delta O$ at the point of the density; along $P_b$, $M_\omega$, $R$ the sector parts in the sector of $A|_{W_{\mathbb R}^\perp}$, with $\{S^{(1)}(M_\omega),O\}=(\omega x)^\mu\partial_\mu O$ on a scalar bilinear and the position part $(\omega x_O)^\mu\Upsilon_{P_{e_\mu}}$ of $\Upsilon_{M_\omega}$; along $K_b$, the part linear in $x$ of Proposition \ref{prop:c11-oddsol} (v) gives $\tfrac{44}3(b,x)O$, which with $-(b,x)\Upsilon^0_\Delta$ and $v(x)^\mu\Upsilon_{P_{e_\mu}}$ is $-\tfrac13\{S^{(1)}(K_b),O\}$ on the traceless part and the same plus $48(b,x)O$ on the trace, $\{S^{(1)}(K_b),O\}=v^\mu\partial_\mu O-4(b,x)O$. (i) Theorem \ref{thm:c9-D1} and Theorem \ref{thm:c11-oddar2}. (ii) $c_Q(\xi_+)+c_Q(\xi_-)-p_1=\tfrac{79}6-6L+\tfrac13-(t_{\xi_+}+t_{\xi_-}-t_P)$ and $c_Q(\xi_+)+c_S(\eta_+)-\tfrac14D_1=\tfrac{79}6-6L+\tfrac13-(t_{\xi_+}+t_{\eta_+}-t_\Delta)$, both $\tfrac{27}2-6L-\tfrac{27}4-2L$ by \eqref{eq:c11-trel}. (iii) Theorem \ref{thm:c11-oddar2} (i), (iii). (iv) Theorem \ref{thm:c11-oddar2} (ii), (iv) and Proposition \ref{prop:c11-freedom}.
\end{proof}
